\SourceLeaf{028}
\textit{superficiei secundae} a priori distinguam. Tunc manifestum est, totum negotium in eo versari, ut demonstretur, ad minimum unum punctum dari, quod simul in plano, in superficie prima et in superficie secunda iaceat.

\Article{17}
Facile perspici potest, superficiem primam partim supra planum partim infra planum iacere; patet enim distantiam a centro $r$ tam magnam accipi posse, ut reliqui termini in $T$ prae primo $r^m\sin m\varphi$ evanescant; hic vero, angulo $\varphi$ rite determinato, tam positivus quam negativus fieri potest. Quare planum fixum necessario a superficie prima secabitur; hanc plani cum superficie prima intersectionem vocabo \textit{lineam primam}; quae itaque determinabitur per aequationem $T=0$. Ex eadem ratione planum a superficie secunda secabitur; intersectio constituet curvam per aequationem $U=0$ determinatam, quam \textit{lineam secundam} appellabo. Proprie utraque curva ex pluribus ramis constabit, qui omnino seiuncti esse possunt, singuli vero erunt lineae continuae. Quin adeo linea prima semper erit talis, quam complexam vocant, axisque $GC$ tamquam pars huius curvae spectanda; quicunque enim valor ipsi $r$ tribuatur, $T$ semper fiet $=0$, quando $\varphi$ aut $=0$ aut $=180^{\circ}$. Sed praestat complexum cunctorum ramorum per omnia puncta, ubi $T=0$, transeuntium tamquam unam curvam considerare (secundum usum in geometria sublimiori generaliter receptum), similiterque cunctos ramos per omnia puncta transeuntes, ubi $U=0$. Patet iam, rem eo reductam esse, ut demonstretur, ad minimum unum punctum in plano dari, ubi ramus aliquis lineae primae a ramo lineae secundae secetur. Ad hunc finem indolem harum linearum propius contemplari oportebit.

\Article{18}
Ante omnia observo, utramque curvam esse algebraicam, et quidem, si ad coordinatas orthogonales revocetur, ordinis $m^{\mathrm{ti}}$. Sumto enim initio abscissarum in $C$, abscissisque $x$ versus $G$, applicatis $y$ versus $P$, erit $x=r\cos\varphi$, $y=r\sin\varphi$, adeoque generaliter, quidquid sit $n$,
\[
\begin{aligned}
r^n\sin n\varphi={}&nx^{n-1}y-\frac{n\mathbin{.}n-1\mathbin{.}n-2}{1\mathbin{.}2\mathbin{.}3}x^{n-3}y^3\\
&+\frac{n\mathbin{.}\mathbin{.}\mathbin{.}n-4}{1\mathbin{.}\mathbin{.}\mathbin{.}\mathbin{.}5}x^{n-5}y^5-\text{etc.},\\
r^n\cos n\varphi={}&x^n-\frac{n\mathbin{.}n-1}{1\mathbin{.}2}x^{n-2}yy\\
&+\frac{n\mathbin{.}n-1\mathbin{.}n-2\mathbin{.}n-3}{1\mathbin{.}2\mathbin{.}3\mathbin{.}4}x^{n-4}y^4-\text{etc.}
\end{aligned}
\]
Quamobrem tum $T$ tum $U$ constabunt ex pluribus huiusmodi terminis $ax^{\alpha}y^{\beta}$,
\SourceLeaf{029}
denotantibus $\alpha$, $\beta$ numeros integros positivos, quorum summa, ubi maxima est, fit $=m$. Ceterum facile praevideri potest, cunctos terminos ipsius $T$ factorem $y$ involvere, adeoque lineam primam proprie ex recta (cuius aequatio $y=0$) et curva ordinis $m-1^{\mathrm{ti}}$ compositam esse; sed necesse non est ad hanc distinctionem hic respicere.

Maioris momenti erit investigatio, an linea prima et secunda crura infinita habeant, et quot qualiaque. In distantia infinita a puncto $C$ linea prima, cuius aequatio $\sin m\varphi+\frac{A}{r}\sin(m-1)\varphi+\frac{B}{rr}\sin(m-2)\varphi\ \text{etc.}=0$, confundetur cum linea, cuius aequatio $\sin m\varphi=0$. Haec vero exhibet $m$ lineas rectas in puncto $C$ se secantes, quarum prima est axis $GCG'$, reliquae contra hanc sub angulis $\frac{1}{m}180$, $\frac{2}{m}180$, $\frac{3}{m}180$ etc. graduum inclinatae. Quare linea prima $2m$ ramos infinitos habet, qui peripheriam circuli radio infinito descripti in $2m$ partes aequales dispertiuntur, ita ut peripheria a ramo primo secetur in concursu circuli et axis, a secundo in distantia $\frac{1}{m}180^{\circ}$, a tertio in distantia $\frac{2}{m}180^{\circ}$ etc. Eodem modo linea secunda in distantia infinita a centro habebit asymptotam per aequationem $\cos m\varphi=0$ expressam, quae est complexus $m$ rectarum in puncto $C$ sub aequalibus angulis itidem se secantium, ita tamen, ut prima cum axe $CG$ constituat angulum $\frac{1}{m}90^{\circ}$, secunda angulum $\frac{3}{m}90^{\circ}$, tertia angulum $\frac{5}{m}90^{\circ}$ etc. Quare linea secunda etiam $2m$ ramos infinitos habebit, quorum singuli medium locum inter binos ramos proximos lineae primae occupabunt, ita ut peripheriam circuli radio infinite magno descripti in punctis, quae $\frac{1}{m}90^{\circ}$, $\frac{3}{m}90^{\circ}$, $\frac{5}{m}90^{\circ}$ etc. ab axe distant, secent. Ceterum palam est, axem ipsum semper duos ramos infinitos lineae primae constituere, puta primum et $m+1^{\mathrm{tum}}$. Luculentissime hic ramorum situs exhibetur in fig. 2, pro casu $m=4$ constructa, ubi rami lineae secundae, ut a ramis lineae primae distinguantur, punctati exprimuntur, quod etiam de figura quarta est tenendum\HistoricalNote{8}. --- Quum vero hae conclusiones maximi momenti sint, quantitatesque infinite magnae quosdam lectores offendere possint: illas etiam absque infinitorum subsidio in art. sequ. eruere docebo.

\Article{19}
\textsc{Theorema.} \textit{Manentibus cunctis ut supra, ex centro $C$ describi poterit circulus,}
